% Ad Atticum 13.38 --- headnote
% ad-atticum-13-38, 4 August 45 BC, Tusculum

\textit{Cicero to Atticus, written at the Tusculan villa around 4
August 45 BC --- Perseus dateline \textit{Scr. in Tusculano
circ. prid. Non. Sext. a. 709 (45)}. Two sections, on two
subjects --- one literary, one domestic --- both written before
breakfast. Cicero is up before dawn working on the
\textit{Academica} (``writing against the Epicureans''); a
courier brings an insulting letter from his nephew, young
Quintus, whose dinner-party rage Cicero had reported two days
earlier in 13.37. Cicero copies the letter on to Atticus and
quotes its choicest line, an unfinished bit of insolence: ``For
my part, whatever ill thing can be said against you --- ''. He
calls it \textit{impurius} --- filthier than anything.}

\textit{The second section is the more revealing. Cicero asks
Atticus to choose his tactics for him --- open contempt or
\textit{crooked stratagems} (the Pindaric phrase from
\textit{Nemean} 7) --- and confesses what he is really afraid of:
being caught alone at the Tusculan villa when the nephew arrives,
or worse, when Caesar himself drops in unannounced on a return
march from Spain. The Greek does the work the Latin can't:
\textit{poteron} for the alternatives, \textit{skoliais apatais}
for the crooked path, the Pindar tag for the divided mind.
The register is hurried, anxious, almost diaristic.}
